\subsection{Recovery Manager: Preassigned Takeover}
\label{sec:parallel-recovery}
\label{sec:recovery-manager}

Recovery Manager (2) decides where each request will resume before
failure occurs. When request $r$ arrives and the Router selects source
DP rank $g_r$, the Manager records a target assignment $p_r\ne g_r$
in control state outside the source's failure domain. The same $p_r$
is both the retry destination and the DP rank that will hold any
Peer-HBM backup of $r$. A simple mapping for $n\ge2$ DPs is
$p_r=(g_r+1)\bmod n$; per-request assignments may distribute load across
several targets. The assignment does not reserve HBM: the KV Master
chooses whether to request Peer protection, and the target Adapter
controls its capacity. A TP group restores compatible shards as one DP.

\paragraph{Activate the recorded route after failure.}
The Manager combines membership, device errors, and GPU progress
reported by Adapters; a process heartbeat alone does not prove GPU
progress. After confirming and isolating a failed source, it obtains
request token prefixes, commit frontiers, model/layout IDs, and replica
locations from surviving control state, then submits retries to their
recorded targets. Available Peer replicas are therefore local to takeover.
DRAM or recomputation supplies the remaining state. Control-state
survival, fencing, failure-domain separation, and retry deduplication
are runtime contracts. A failed target requires a new valid assignment;
copies on the old target cannot be treated as local to the new target.

Retries enter independently, subject to foreground admission, continuous
batching, and shared compute/copy capacity. Each request advances on its
own layer-readiness events. Slack determines protection, not a request-level
completion barrier; the Manager never waits for earlier retries merely
because they had higher placement priority.
